\section{模型总结与评价}

\subsection{模型总结}

本文按照“先识别数据结构，再逐步释放决策自由度”的思路构造四问模型链。Q1 从任务到达、区域--类型组成和条件资源分布中提取短期需求结构，并以标记复合 Poisson 给出 24\,h 概率预测和基础算力调度；Q2 在储能动作固定的条件下释放任务跨区与开工时段自由度，以动态边际能源响应修复附件中显著的“弃电与购电并存”时空错配；Q3 固定 workload 后单独释放新能源分配与 BESS 时间柔性，并由 $H_{rt}=R_{rt}-L_{rt}$ 与 SellLimit 的数据结构自然识别 A/B/C、D/E、F 三类储能角色；Q4 最后同时释放 workload 与 BESS/Grid 决策，在统一的计算--能源 recourse 中采用 region multi-cut Benders 与 exact column generation 求解，并以 Cost$\rightarrow$Wait$\rightarrow$Latency 字典序完成服务质量精化。

Q4 的 50,000 任务 canonical 联合端点经两轮 final-pool 再认证后保持一致：真实能源成本为 $-4.5934\times10^8$ CNY，Total Wait 为 0，仅 61 个任务跨区执行，迁移率为 0.122\%，平均网络时延为 5.0172 ms。与 Q2 在冻结能源动作时约 74.896\% 的迁移率和 20.475 h 的平均等待相比，该结果表明本附件数据下\textbf{储能时间柔性替代了绝大部分原本由 workload 时间平移和跨区迁移承担的调节需求}；因此 Q2 与 Q3 的单侧价值不能简单相加，联合优化改变的是两类柔性的分工。

\subsection{模型优点}

\begin{enumerate}[label=\textbf{\arabic*.}]
\item \textbf{数据结构直接决定模型，而不是先选算法再找解释。} Q1 的 Region$\times$TaskType 关联和任务数--GPU-hour 错位决定联合标记与资源 mark；Q2 的 91.9\% region-hour “弃电+购电”并存以及弃电的低碳属性决定动态边际响应；Q3 的 $H$ 与 SellLimit 直接导出储能角色；Q4 的区域优势错位和顺序 Q2$\rightarrow$Q3 失真探针决定必须联合优化。

\item \textbf{避免主观多指标权重。} Cost、Carbon、Wait、Latency、RenewableUtilization 与 PeakNetImport 不通过 AHP、熵权或等权综合分强行揉成单目标；Cost/Wait/Latency 使用字典序，Carbon 在冲突被数据或场景激活时使用 epsilon-constraint，其余指标统一报告并按辨识力决定是否进入约束层。

\item \textbf{完整合法域不作经验裁剪。} Q2 保留完整合法区域与合法开始时刻；Q4 面对约 $2.33\times10^8$ 个 placements，不采用 top-k 或任意等待窗，而以 exact pricing 隐式搜索完整域。Facility/IT 重复约束、稀疏 Benders dual support 等压缩均来自严格数据结构，不改变可行域。

\item \textbf{验证链完整且结果边界清楚。} Q1 使用独立验证/测试与概率区间校准；Q2 使用多起点与完整扫描收敛门禁；Q3 用 B0/E0/E1 分离 raw reference、无储能优化和储能增量价值；Q4 使用 40-task exact benchmark、真实 Energy LP 复核、region multi-cut、两轮 final-pool 再认证和全硬约束审计。

\item \textbf{联合模型给出了非显然的机制结论。} Q4 并未继续扩大 workload 迁移和等待，而是在同一成本面内恢复到 Wait=0、migration=0.122\%。这一结果把“算力柔性”和“储能柔性”的替代关系从概念推断变成了任务级排程与能源 recourse 共同支持的数值结论。
\end{enumerate}

\subsection{模型局限}

\begin{enumerate}[label=\textbf{\arabic*.}]
\item \textbf{Q1 的最终独立测试窗口较短。} 测试集仅包含 2376--2399 共 24 个小时，System 层 90\% 区间覆盖率为 87.5\%（21/24）。该样本量足以作为题目要求的短期 baseline 检验，但不足以支持长期季节性或更复杂到达机制的结论。

\item \textbf{Q2 的构造式端点存在路径依赖。} reference-start 与 previous-start 的方向性降本减碳结论一致，但精确整数端点存在约 1\% 量级差异，因此 Q2 的结果应解释为通过完整合法域动态边际扫描得到的高质量可行解，而非全局整数最优证书。

\item \textbf{Q4 的 50k 全局整数最优尚未证明。} 两轮 final-pool 再认证中完整域 LP pricing 与 region Benders 均闭合，restricted-column MIP gap 为 0；但 Latency 的完整域 LP 下界为 250469.9546 ms，整数代表解为 250860 ms，相对差约 0.1555\%。因此本文只称其为 final-pool recertified best-known integer representative solution，不将 restricted MIP gap=0 等同于完整域整数最优。

\item \textbf{跨问绝对成本仍需统一 Hour 2406 accounting。} Q3 E1 与 Q4 canonical 的表面成本差很小，但在 2406 terminal settlement、购售电统计时域和 Baseline 命名完全对齐前，不应把该差值直接解释为联合优化的正式增量收益。

\item \textbf{场景结论依赖透明的题内构造。} Q4 canonical 主问题已经完成，但 Carbon constraints、电价机制和新能源波动三类正式场景仍需在同一 joint model 中重新优化后才能冻结；历史 $\gamma=1.4$ 只承担诊断作用，不能替代正式场景。
\end{enumerate}

\subsection{改进方向}

若进一步扩展，可在保持现有完整合法域和真实 Energy recourse 的前提下增加更长时间尺度的数据验证、采用更多独立样本检验 Q1 的到达稳定性，并在 Q4 正式场景完成后分析柔性替代关系在不同碳预算、电价机制和新能源波动下是否保持。对于 50k 整数最优性，只有在需要更严格理论证书且计算预算允许时，才有必要进一步研究 branch-price-and-Benders；当前约 0.156\% 的 Latency LP--integer gap 已足以支撑比赛论文中的谨慎最优性边界说明。